\documentclass[10pt]{beamer}
\usetheme{metropolis}
\usepackage{graphicx}
\usepackage{listings}
\usepackage{booktabs}
\usepackage{xcolor}
\definecolor{codebg}{HTML}{F5F5F5}
\definecolor{codegreen}{HTML}{2E7D32}
\definecolor{codegray}{HTML}{757575}
\definecolor{codeblue}{HTML}{1565C0}
\lstdefinestyle{rstyle}{language=R, backgroundcolor=\color{codebg}, basicstyle=\ttfamily\scriptsize,
  keywordstyle=\color{codeblue}\bfseries, stringstyle=\color{codegreen},
  commentstyle=\color{codegray}\itshape, breaklines=true, frame=single,
  rulecolor=\color{codegray!30}, showstringspaces=false, tabsize=2,
  morekeywords={library,ggplot,aes,geom_col,geom_bar,geom_segment,geom_point,geom_text,
    coord_flip,fct_reorder,scale_fill_manual,position_dodge,position_stack,position_fill,
    theme_minimal,theme,element_text,element_blank,labs,count,slice_max,mutate}}
\lstdefinestyle{pythonstyle}{language=Python, backgroundcolor=\color{codebg}, basicstyle=\ttfamily\scriptsize,
  keywordstyle=\color{codeblue}\bfseries, stringstyle=\color{codegreen},
  commentstyle=\color{codegray}\itshape, breaklines=true, frame=single,
  rulecolor=\color{codegray!30}, showstringspaces=false, tabsize=4,
  morekeywords={import,from,as,pd,plt,sns,np,fig,ax,barh,bar,hlines,scatter,text}}
\title{Comparisons: Bar Charts, Lollipop \& Cleveland Dot}
\subtitle{Workshop 3 --- Module 3}
\author{Prof.Asoc.\ Endri Raco}
\institute{Data Visualization for Data Scientists \\ UNYT Tirana}
\date{2025/26}
\begin{document}

\begin{frame}
  \titlepage
\end{frame}

\begin{frame}{Module Roadmap}
  \begin{enumerate}
    \item Bar chart best practices: horizontal, sorted, labelled
    \item Three bar variants: grouped, stacked, 100\% stacked
    \item Lollipop charts: less ink, same precision
    \item Cleveland dot plots: Tufte-minimal ranking
    \item Dumbbell charts: before/after comparisons
    \item Comparison chart selection guide
  \end{enumerate}
  \begin{exampleblock}{Learning Outcome}
    You will select and produce the optimal comparison chart for any
    categorical variable, from the standard bar chart to the Tufte-minimal
    Cleveland dot plot.
  \end{exampleblock}
\end{frame}

\section{Bar Chart Best Practices}

\begin{frame}{Vertical vs Horizontal}
  \begin{center}
    \includegraphics[width=0.92\textwidth]{figures_w03m03/vert_vs_horiz}
  \end{center}
  \begin{alertblock}{Pro-Tip}
    \textbf{Always use horizontal bars} for categorical data with text
    labels. Vertical bars require rotated labels (45\textdegree{} or 90\textdegree{}),
    which are harder to read. Use vertical only when x is temporal
    (years, months) or ordinal.
  \end{alertblock}
\end{frame}

\begin{frame}{Sorted vs Unsorted}
  \begin{center}
    \includegraphics[width=0.92\textwidth]{figures_w03m03/sorted_vs_unsorted}
  \end{center}
  \begin{exampleblock}{Data Insight}
    Alphabetical order is almost never meaningful. \textbf{Sort by value}
    to make ranking instantly visible. Use \texttt{fct\_reorder(cat, n)}
    in R or \texttt{sort\_values()} in Python. Add \textbf{grey-accent}
    colour to highlight the key finding.
  \end{exampleblock}
\end{frame}

\begin{frame}[fragile]{Sorted Horizontal Bar in R vs Python}
  \begin{columns}[T]
    \begin{column}{0.48\textwidth}
      \textbf{R (ggplot2)}
\begin{lstlisting}[style=rstyle]
apps |>
  count(Category, sort = TRUE) |>
  slice_max(n, n = 10) |>
  mutate(
    Category = fct_reorder(
      Category, n),
    hl = n == max(n)) |>
  ggplot(aes(x = Category, y = n,
             fill = hl)) +
  geom_col(width = 0.5,
           show.legend = FALSE) +
  geom_text(aes(label =
    scales::comma(n)),
    hjust = -0.15, size = 3,
    fontface = "bold") +
  scale_fill_manual(values =
    c("TRUE"="#1565C0",
      "FALSE"="#BBBBBB")) +
  coord_flip() +
  theme_minimal() +
  theme(panel.grid.major.y =
    element_blank())
\end{lstlisting}
    \end{column}
    \begin{column}{0.48\textwidth}
      \textbf{Python}
\begin{lstlisting}[style=pythonstyle]
top10 = (apps
    .value_counts("Category")
    .head(10).sort_values())

fig, ax = plt.subplots()
colors = ["#BBBBBB"] * 10
colors[-1] = "#1565C0"

ax.barh(top10.index, top10.values,
    color=colors, height=0.55)

for i, v in enumerate(
    top10.values):
    ax.text(v + 15, i, f"{v:,}",
        va="center", fontsize=7,
        fontweight="bold",
        color="#1565C0"
          if i == 9 else "#888")

for sp in ["top","right","left"]:
    ax.spines[sp].set_visible(False)
ax.set_xticks([])
ax.tick_params(left=False)
\end{lstlisting}
    \end{column}
  \end{columns}
\end{frame}

\section{Bar Chart Variants}

\begin{frame}{Grouped, Stacked, and 100\% Stacked}
  \begin{center}
    \includegraphics[width=0.95\textwidth]{figures_w03m03/grouped_stacked}
  \end{center}
\end{frame}

\begin{frame}{When to Use Each Variant}
  \centering
  \begin{tabular}{lp{4cm}p{4cm}}
    \toprule
    \textbf{Variant} & \textbf{Question} & \textbf{Limitation} \\
    \midrule
    Grouped (dodge) & Compare values between groups within each category & Max 3 groups; more = clutter \\
    Stacked & Compare totals across categories + see group contribution & Hard to compare non-baseline segments \\
    100\% Stacked & Compare proportions & Loses magnitude (totals invisible) \\
    \bottomrule
  \end{tabular}

  \vspace{6pt}
  \begin{alertblock}{Pro-Tip}
    If you have only \textbf{two groups}, use a grouped bar --- proportions
    are visible without 100\% stacking. If proportions are the message,
    prefer 100\% stacked. If totals matter, use regular stacked.
  \end{alertblock}
\end{frame}

\begin{frame}[fragile]{Bar Variants in R vs Python}
  \begin{columns}[T]
    \begin{column}{0.48\textwidth}
\begin{lstlisting}[style=rstyle]
# Grouped
ggplot(df, aes(x = cat,
    fill = group)) +
  geom_bar(position = "dodge")

# Stacked
ggplot(df, aes(x = cat,
    fill = group)) +
  geom_bar(position = "stack")

# 100% Stacked
ggplot(df, aes(x = cat,
    fill = group)) +
  geom_bar(position = "fill") +
  scale_y_continuous(
    labels = scales::percent)
\end{lstlisting}
    \end{column}
    \begin{column}{0.48\textwidth}
\begin{lstlisting}[style=pythonstyle]
# Grouped (manual offset)
x = np.arange(len(cats))
ax.bar(x - w/2, free, w,
    label="Free")
ax.bar(x + w/2, paid, w,
    label="Paid")

# Stacked
ax.bar(cats, free, label="Free")
ax.bar(cats, paid, bottom=free,
    label="Paid")

# 100% Stacked (compute %)
ct = pd.crosstab(
    df["cat"], df["group"],
    normalize="index") * 100
ct.plot.bar(stacked=True, ax=ax)
\end{lstlisting}
    \end{column}
  \end{columns}
\end{frame}

\section{Lollipop Charts}

\begin{frame}{Lollipop: Less Ink, Same Precision}
  \begin{center}
    \includegraphics[width=0.78\textwidth]{figures_w03m03/lollipop}
  \end{center}
  \begin{exampleblock}{Data Insight}
    The lollipop chart replaces the bar's filled rectangle with a thin
    line (segment) and a point. This reduces data ink while preserving
    the position encoding. It is especially effective for long category
    lists (10--20 items) where bars become visually heavy.
  \end{exampleblock}
\end{frame}

\begin{frame}[fragile]{Lollipop in R vs Python}
  \begin{columns}[T]
    \begin{column}{0.48\textwidth}
\begin{lstlisting}[style=rstyle]
ggplot(df, aes(
    x = fct_reorder(cat, val),
    y = val)) +
  geom_segment(aes(
    xend = cat, y = 0,
    yend = val),
    color = "grey70",
    linewidth = 0.8) +
  geom_point(size = 3,
    color = "#1565C0") +
  coord_flip() +
  theme_minimal() +
  theme(
    panel.grid.major.y =
      element_blank()) +
  labs(x = NULL, y = "Count")
\end{lstlisting}
    \end{column}
    \begin{column}{0.48\textwidth}
\begin{lstlisting}[style=pythonstyle]
fig, ax = plt.subplots()

# Segments (stems)
ax.hlines(y=range(len(cats)),
    xmin=0, xmax=vals,
    color="#BBBBBB", lw=1.5)

# Points (dots)
ax.scatter(vals,
    range(len(cats)),
    s=60, c="#1565C0",
    zorder=5,
    edgecolors="white", lw=0.5)

# Direct labels
for i, v in enumerate(vals):
    ax.text(v + 10, i, f"{v:,}",
        va="center", fontsize=7,
        fontweight="bold")

ax.set_yticks(range(len(cats)))
ax.set_yticklabels(cats)
\end{lstlisting}
    \end{column}
  \end{columns}
\end{frame}

\section{Cleveland Dot Plot}

\begin{frame}{Cleveland Dot: Minimal Ink, Maximum Precision}
  \begin{center}
    \includegraphics[width=0.78\textwidth]{figures_w03m03/cleveland_dot}
  \end{center}
  \begin{alertblock}{Pro-Tip}
    The Cleveland dot plot uses only a point per category --- the
    absolute minimum of data ink (Tufte, 1983). Add subtle gridlines
    for alignment and a mean reference line for context. Best for
    audiences who value precision over visual impact.
  \end{alertblock}
\end{frame}

\begin{frame}{Bar vs Lollipop vs Cleveland: Ink Reduction}
  \begin{center}
    \includegraphics[width=0.95\textwidth]{figures_w03m03/three_comparison}
  \end{center}
\end{frame}

\section{Dumbbell Charts}

\begin{frame}{Dumbbell: Before/After Comparison}
  \begin{center}
    \includegraphics[width=0.82\textwidth]{figures_w03m03/dumbbell}
  \end{center}
  \begin{exampleblock}{Data Insight}
    Dumbbell charts show change between two time points per category.
    The connector length encodes the magnitude of change; colour
    distinguishes the two periods. Direct-label the change value
    (\texttt{+459}) to avoid requiring mental arithmetic.
  \end{exampleblock}
\end{frame}

\begin{frame}[fragile]{Dumbbell in R vs Python}
  \begin{columns}[T]
    \begin{column}{0.48\textwidth}
\begin{lstlisting}[style=rstyle]
ggplot(df, aes(
    y = fct_reorder(cat,
      val_2024))) +
  geom_segment(aes(
    x = val_2020, xend = val_2024,
    yend = cat),
    color = "grey80",
    linewidth = 1.5) +
  geom_point(aes(x = val_2020),
    size = 3, color = "#1565C0") +
  geom_point(aes(x = val_2024),
    size = 3, color = "#E53935") +
  geom_text(aes(x = val_2024,
    label = paste0("+",
      val_2024 - val_2020)),
    hjust = -0.3, size = 3,
    color = "#E53935") +
  theme_minimal() +
  labs(x = "Count", y = NULL,
    title = "Change: 2020 vs 2024")
\end{lstlisting}
    \end{column}
    \begin{column}{0.48\textwidth}
\begin{lstlisting}[style=pythonstyle]
for i in range(len(cats)):
    # Connector
    ax.plot([v2020[i], v2024[i]],
        [i, i], color="#CCCCCC",
        lw=2.5,
        solid_capstyle="round")
    # Start point (2020)
    ax.scatter([v2020[i]], [i],
        s=60, c="#1565C0", zorder=5,
        edgecolors="white", lw=0.5)
    # End point (2024)
    ax.scatter([v2024[i]], [i],
        s=60, c="#E53935", zorder=5,
        edgecolors="white", lw=0.5)
    # Change label
    ax.text(v2024[i] + 20, i,
        f"+{v2024[i]-v2020[i]:,}",
        va="center", fontsize=7,
        fontweight="bold",
        color="#E53935")
\end{lstlisting}
    \end{column}
  \end{columns}
\end{frame}

\section{Chart Selection}

\begin{frame}{Comparison Chart Selection Guide}
  \begin{center}
    \includegraphics[width=0.92\textwidth]{figures_w03m03/decision_card}
  \end{center}
\end{frame}

\section{Complete Examples}

\begin{frame}[fragile]{R: Publication-Quality Lollipop}
  \begin{columns}[T]
    \begin{column}{0.52\textwidth}
\begin{lstlisting}[style=rstyle]
apps |>
  count(Category, sort = TRUE) |>
  slice_max(n, n = 8) |>
  mutate(Category =
    fct_reorder(Category, n),
    hl = n == max(n)) |>
  ggplot(aes(x = Category,
             y = n)) +
  geom_segment(aes(xend = Category,
    y = 0, yend = n),
    color = "grey75") +
  geom_point(aes(color = hl),
    size = 3, show.legend = FALSE) +
  scale_color_manual(values =
    c("TRUE"="#1565C0",
      "FALSE"="#BBBBBB")) +
  geom_text(aes(
    label = scales::comma(n)),
    hjust = -0.3, size = 3,
    fontface = "bold") +
  coord_flip() +
  theme_minimal() +
  labs(title = "FAMILY Leads",
       x = NULL, y = NULL)
\end{lstlisting}
    \end{column}
    \begin{column}{0.45\textwidth}
      \includegraphics[width=\textwidth]{figures_w03m03/r_output}
    \end{column}
  \end{columns}
\end{frame}

\begin{frame}[fragile]{Python: Same Lollipop}
  \begin{columns}[T]
    \begin{column}{0.52\textwidth}
\begin{lstlisting}[style=pythonstyle]
top8 = (apps
    .value_counts("Category")
    .head(8).sort_values())

fig, ax = plt.subplots(
    figsize=(6, 5))
colors = ["#BBBBBB"] * 8
colors[-1] = "#2E7D32"

ax.hlines(y=range(8), xmin=0,
    xmax=top8.values,
    color="#BBBBBB", lw=1.5)
ax.scatter(top8.values,
    range(8), s=50, c=colors,
    zorder=5, edgecolors="white")

for i, v in enumerate(
    top8.values):
    ax.text(v + 15, i, f"{v:,}",
        va="center", fontsize=7,
        fontweight="bold",
        color="#2E7D32"
          if i == 7 else "#888")

ax.set_yticks(range(8))
ax.set_yticklabels(top8.index)
ax.set_title("FAMILY Leads",
    fontweight="bold")
\end{lstlisting}
    \end{column}
    \begin{column}{0.45\textwidth}
      \includegraphics[width=\textwidth]{figures_w03m03/py_output}
    \end{column}
  \end{columns}
\end{frame}

\section{Synthesis}

\begin{frame}{Key Takeaways}
  \begin{enumerate}
    \item \textbf{Always horizontal + sorted + labelled} for categorical
          bar charts. Vertical only for temporal x-axes.
    \item Three bar variants serve different questions: \textbf{grouped}
          (compare values), \textbf{stacked} (compare totals),
          \textbf{100\% stacked} (compare proportions).
    \item \textbf{Lollipop} = segment + point: less ink than bars,
          same precision. Best for 10--20 categories.
    \item \textbf{Cleveland dot} = point only + gridlines: minimum ink,
          Tufte-approved. Best for precision-oriented audiences.
    \item \textbf{Dumbbell} = two points + connector: shows change
          between two time periods per category.
    \item The progression bar $\rightarrow$ lollipop $\rightarrow$ dot
          is a data-ink reduction cascade.
  \end{enumerate}
\end{frame}

\begin{frame}{Essential Readings}
  \begin{itemize}
    \item Cleveland, W.~S. (1985). \emph{The Elements of Graphing Data}.
          Wadsworth. (Origin of the dot plot.)
    \item Schwabish, J. (2021). \emph{Better Data Visualizations},
          Chapter 4 (Bar Charts and Alternatives).
    \item Wickham, H. (2016). \emph{ggplot2}, Chapter 3 (Bar Charts).
    \item Wilke, C.~O. (2019). \emph{Fundamentals of Data Visualization},
          Chapter 6 (Amounts).
  \end{itemize}
  \vspace{6pt}
  \textbf{Next Module}: Relationships --- Scatter, Bubble \& Hexbin (W03-M04)
\end{frame}

\end{document}
